I cifrari a blocchi fanno uno di una sequenza di \emph{permutazioni} e \emph{sostituzioni}. inoltre spesso sono implementati sotto forma di cifrari iterati. Un cifrario iterato lavora su $\mathcal{N}$ ``round'', per cui è necessario definire una \emph{funzione round} e una \emph{key-schedule} (una tupla delle chiavi relative ai round). La funzione del round $g$ ha due input, la chiave relativa al round $K^r$ e uno \emph{stato corrente} $w^{r-1}$. Inoltre per far si che il messaggio si possa decifrare $g$ deve essere iniettiva rispetto alla chiave.

\subsection{Sostitution Permutation Network}
Il cifrario \emph{SPN} è una modifica di un cifrario iterato in cui i \emph{plaintext} e \emph{ciphertext} sono divisi in $m$ blocchi lunghi $\ell$.
\begin{definition}[SPN]
    Siano $\ell,m, \mathcal{N}$ interi positivi $M = C = \{0,1\}^{\ell m}$, siano $\pi_S, \pi_P$ funzioni di permutazione
    \begin{align*}
        \pi_S: \{0,1\}^{\ell} &\to \{0,1\}^{\ell}\\
        \pi_P: \{1,\dots,\ell m\} &\to \{1,\dots, \ell m\}
    \end{align*}
    con $\mathcal{K} \subseteq (\{0,1\}^{\ell m})^{\mathcal{N}+1}$ i \emph{key-schedule} possibili derivati da una chiave iniziale K.

    La funzione di cifratura è descritta dall'algoritmo \ref{alg:spn}, e il l'\textit{i-esimo} blocco è descritto come $x_{<i>}$ per $i = 1,\dots,m$
\end{definition}
\begin{figure}[ht]
    \begin{algorithm}[H]
        \label{alg:spn}
        
        \DontPrintSemicolon
        \TitleOfAlgo{SPN($x, \pi_S, \pi_P, (K^1, \dots, K^{\mathcal{N}+1})$)}
        
        $w^0 \leftarrow x$\;
        \For{$r \leftarrow 1$ \KwTo $\mathcal{N}-1$}{
            $u^r \leftarrow w^{r-1} \oplus K^r$\;
            \For{$i \leftarrow 1$ \KwTo $m$}{
                $v^r_{\langle i \rangle} \leftarrow \pi_S(u^r_{\langle i \rangle})$\;
                }
                $w^r \leftarrow (v^r_{\pi_P(1)}, \dots, v^r_{\pi_P(\ell m)})$\;
                }
                $u^{\mathcal{N}} \leftarrow w^{\mathcal{N}-1} \oplus K^{\mathcal{N}}$\;
                \For{$i \leftarrow 1$ \KwTo $m$}{
                    $v^{\mathcal{N}}_{\langle i \rangle} \leftarrow \pi_S(u^{\mathcal{N}}_{\langle i \rangle})$\;
                    }
                    $y \leftarrow v^{\mathcal{N}} \oplus K^{\mathcal{N}+1}$\;
                    \Return{$y$}\;
                \end{algorithm}
\end{figure}